\section{Ресурсно-нормоване вимірювання та двофакторний закон}
\label{sec:model}

Введено протокол ізоляції, двофакторну характеристику та результат
ідентифікованості за сформульованих припущень.

\subsection{Ресурсно-нормоване вимірювання як інструмент ізоляції}

Якщо $\n$ валідаторів ділять хост із $Q$ ядрами, кожен отримує $Q/\n$ ядер, тож
виміряна пропускна здатність змішує алгоритмічну вартість зі скорочуваним
обчислювальним бюджетом. \RNM{} призначений ізолювати ці ефекти за контрольованих
квот.

\begin{Def}[Ресурсно-нормоване вимірювання]
\label{def:rnm}
Кампанія відповідає \emph{протоколу \RNM{}} з квотою~$\q$, якщо кожен контейнер
валідатора отримує однакову CPU-квоту та ліміт пам'яті незалежно від~$\n$, а
сумарна резервація $\n\q$ не перевищує ємності хоста. Нормована пропускна
здатність --- $\TPSn\eqdef\TPS/\q$.
\end{Def}

\begin{Assumption}[Сепарабельний обчислювальний бюджет]
\label{as:separable}
На робочому діапазоні квот у кампанії
$\TPS(\n,\cc,\q)=\q\cdot\Phi(\n,\cc)\cdot(1+r)$ з $|r|\le C\q$ для сталої~$C$,
незалежної від~$(\n,\cc)$ на цьому діапазоні.
\end{Assumption}

\paragraph{Перевірка припущення.}
За Припущенням~\ref{as:separable} і степеневим фактором
$\Phi\propto\n^{-\beta}$ нахил МНК $\log\TPSn$ за $\log\n$ при фіксованій
конкурентності не залежить від~$\q$: відносний залишок поглинається вільним
членом. Розділ~\ref{sec:exp} перевіряє цю інваріантність для
$\q\in\{0.20,0.25,0.33\}$ і відхиляє її при $\q=0.33$: \RNM{} стабілізує
$\beta$ за фіксованої квоти, а не робить його незалежним від квоти. Калібровану
модель звітовано при $\q=0.25$; поза цим режимом коефіцієнти потрібно
перепідігнати.

\begin{Def}[Двофакторна характеристика масштабування]
\label{def:bifactor}
За протоколом \RNM{}
\begin{equation}
\TPS(\n,\cc)=\Tzero\cdot\gfun(\cc)\cdot\n^{-\beta},
\quad
\gfun(\cc)=\frac{\cc^{\gamma}}{1+(\cc/\csat)^{\theta}},
\label{eq:bifactor}
\end{equation}
де $\Tzero>0$, $\beta>0$, $\gamma\in(0,1]$, $\theta>\gamma$, $\csat>0$.
У ненасиченому режимі $\cc\ll\csat$ залежність лог-лінійна:
$\log\TPS=\log\Tzero+\gamma\log\cc-\beta\log\n$.
\end{Def}

Форма відокремлює конкурентність від масштабування за кількістю валідаторів,
зводиться до степеневого закону при $\cc\ll\csat$, насичується при великому~$\cc$
і містить класичну однофакторну підгонку як випадок $\gamma=0$ (або $\lambda=0$
з $\alpha=-\beta$).

\subsection{Ідентифікованість}

Пропозиція~\ref{prop:identifiability} --- це хрестоматійне зміщення через
пропущену змінну при регресії $\log\TPS$ на $\log\n$ без
$\log\cc$~\cite{wooldridge2010}, чинне за ненасиченої двофакторної моделі та
зазначених припущень щодо шляху й шуму. Її значення інтерпретаційне: вона
показує, як читати виміряний однофакторний показник за заданого шляху
конкурентності. Розділ~\ref{sec:exp} перевіряє узгодженість лише в
дослідженому розгортанні.

\begin{Proposition}[Ідентифікованість показника консенсусу]
\label{prop:identifiability}
Нехай спостереження відповідають ненасиченому двофакторному закону з адитивним
шумом нульового середнього і скінченної дисперсії вздовж шляху конкурентності
$\log\cc_{i}=\lambda\log\n_{i}+\mu+\eta_{i}$
($\E[\eta_{i}\mid\n_{i}]=0$). Якщо емпіричний другий момент
$x_{i}=\log\n_{i}$ збігається до додатної границі, то нахил МНК $\ahat_{m}$
залежності $\log\TPS_{i}$ лише від $x_{i}$ задовольняє
$\ahat_{m}\xrightarrow{\Prob}\gamma\lambda-\beta$.
Зокрема $\ahat>0\iff\lambda>\beta/\gamma$ і $\ahat=-\beta\iff\lambda=0$.
\end{Proposition}

\begin{Proof}
При $x_{i}=\log\n_{i}$ шлях зводить ненасичений закон до
$\log\TPS_{i}=a+\alpha x_{i}+\varepsilon_{i}$, $\E[\varepsilon_{i}\mid x_{i}]=0$,
зі скінченною дисперсією та
$\alpha=\underbrace{-\beta}_{\text{протокол}}
+\underbrace{\gamma\lambda}_{\text{вимірювання}}$.
Спроможність МНК за умови на другий момент дає
$\ahat_{m}\xrightarrow{\Prob}\alpha$~\cite{wooldridge2010}; знаки випливають з
$\gamma>0$.
\end{Proof}

У межах зазначеної моделі лише факторний дизайн ($\lambda=0$) робить
однофакторний нахил спроможною оцінкою~$-\beta$. Для ілюстрації гіпотетичні
значення $\gamma=0.8$, $\beta=0.5$ і $\lambda=1$ (не підігнані) дають
$\ahat=+0.3$, а екстраполяція з $\n=5$ до $\n=60$ відрізняється приблизно в
$7$~разів, що відповідає порядку розбіжності в~\cite{peniak2026a}.

\begin{Corollary}[Непереносність однофакторних показників]
\label{cor:sign}
У ненасиченій двофакторній моделі лабораторії, що вимірюють уздовж
$\lambda_{1}\ne\lambda_{2}$, отримують показники, що відрізняються на
$\gamma(\lambda_{1}-\lambda_{2})$. Без зазначення~$\lambda$ опубліковані
однофакторні показники не слід трактувати як переносні властивості протоколу між
дизайнами вимірювання.
\end{Corollary}

\paragraph{Реалізація.}
На факторній сітці ($\lambda=0$) оцінюємо $(\log\Tzero,\gamma,\beta)$ методом
МНК у лог-просторі та $(\csat,\theta)$ нелінійним МНК на повному проході
за~$\cc$ (розділ~\ref{sec:exp}).
